\subsubsection{Nonword Stimuli}

For nonword stimuli, both OS and LP modulated the N250 complexity 
effect, but through qualitatively distinct mechanisms that reflect the 
different processing profiles of the two dimensions. This constitutes 
a partial qualification of the predicted OS-specific modulation at the 
N250, but the nature of the qualification is itself theoretically 
informative: rather than reflecting a single shared mechanism operating 
at different strengths, the two dimensions produce reversals of the 
complexity effect through entirely distinct trajectories.

For readers lower in OS, complex nonwords elicited substantially more 
negative N250 amplitudes than simple nonwords, reflecting a large early 
processing cost associated with the presence of a real suffix. As OS 
increased, this complexity cost attenuated monotonically, reversing 
significantly at the upper extreme of the observed OS range (OS $= 
+1.9$ SD: $\Delta = -0.490\,\mu$V, $t = -4.27$, $p < .001$), though 
not at OS $= +1$ SD ($\Delta = -0.142\,\mu$V, $t = -1.93$, $p = 
.053$). Examination of estimated marginal means indicated that this 
reversal was driven primarily by the trajectory of complex nonwords, 
which decreased in negativity across the OS range (EMM: $-1.720\,\mu$V 
at OS $= -1.9$ SD to $+0.160\,\mu$V at OS $= +1.9$ SD), while simple 
nonwords also decreased in negativity but more slowly (EMM: 
$-0.720\,\mu$V at OS $= -1.9$ SD to $-0.340\,\mu$V at OS $= +1.9$ 
SD). We propose that this pattern reflects affix-led parsing 
\citep{taft_lexical_1975, rastleMorphologicalDecompositionBased2008}, 
in which the presence of a real suffix triggers morpho-orthographic 
parsing activity across all readers, but the efficiency with which 
that parsing proceeds varies with OS. For readers lower in OS, the 
parse is triggered by the real suffix but cannot be resolved fluently, 
producing a processing cost for complex nonwords relative to simple 
ones. As OS increases, parsing becomes more efficient, reducing and 
eventually reversing the complexity cost. At very high OS, the parse 
is sufficiently efficient that complex nonwords become easier to 
process than simple ones, which lack a real suffix and therefore do 
not benefit from affix-based parsing in the same way. The reversal is 
therefore a consequence of the complex nonword trajectory outpacing 
the simple nonword trajectory as OS increases --- not of a rising 
processing cost for simple nonwords.

The LP pattern was qualitatively different, and the contrast with the 
OS pattern constitutes a dissociation of mechanism within the nonword 
N250 data. As with OS, the complexity effect reversed at high LP --- 
significantly so at LP $= +1$ SD ($\Delta = -0.141\,\mu$V, $t = 
-2.14$, $p = .032$) and more strongly at LP $= +3$ SD ($\Delta = 
-0.860\,\mu$V, $t = -6.38$, $p < .001$). However, examination of 
estimated marginal means indicated that the mechanism was entirely 
distinct from that underlying the OS reversal. Complex nonwords were 
essentially flat across the LP range (EMM: $-0.690\,\mu$V at LP $= 
-2.5$ SD to $-0.820\,\mu$V at LP $= +3.0$ SD), indicating that LP 
did not modulate the processing of complex nonwords. Instead, the 
reversal was driven entirely by simple nonwords becoming increasingly 
negative as LP increased, from $+0.420\,\mu$V at LP $= -2.5$ SD to 
$-1.700\,\mu$V at LP $= +3.0$ SD. We propose that this pattern 
reflects stem-led parsing: readers higher in LP have richer and more 
precise representations of stems as free morphemes with discrete 
lexical entries, and recognition of the stem in a nonword triggers an 
attempt to assign morphological structure to the whole string 
regardless of what follows. This proposal is consistent with 
\citet{beyersmannEmbeddedStemPriming2016}, who found that high 
proficiency readers showed embedded stem priming effects independent 
of whether the stem was followed by a real affix or a non-affixal 
ending --- a pattern suggesting that stem activation in skilled readers 
is not contingent on affix recognition, in contrast to the standard 
affix-stripping account \citep{taft_lexical_1975, 
rastleMorphologicalDecompositionBased2008}. For complex nonwords, the 
stem-led parse succeeds --- a real stem is followed by a real suffix, 
yielding a valid morphological analysis --- and successful parsing 
produces neither a cost nor a benefit at the N250, keeping amplitude 
stable across LP levels. For simple nonwords, the parse fails --- the 
pseudoaffix cannot be integrated into a valid morphological structure 
--- and this failed parse produces a processing cost that grows with 
LP, driving the simple nonword trajectory increasingly negative.

The contrast between the two reversal mechanisms is theoretically 
precise. The OS reversal is driven by the complex nonword trajectory: 
the real suffix triggers morpho-orthographic parsing in all readers, 
but higher OS readers resolve that parse more efficiently, driving 
complex nonwords to decrease in negativity faster than simple nonwords 
and eventually producing a reversal. Simple nonwords are not 
specifically affected by affix-based parsing efficiency because they 
do not contain a real suffix. The LP reversal is driven by the simple 
nonword trajectory: stem recognition triggers an attempt to assign 
morphological structure to the whole string in higher LP readers, and 
this attempt succeeds for complex nonwords --- keeping their amplitude 
stable across LP levels --- but fails for simple nonwords, producing 
a processing cost that grows with LP. This dissociation maps directly 
onto the theoretical profiles of the two dimensions: OS indexes the 
efficiency of sublexical orthographic processing, and its effect on 
early parsing is expressed specifically in the complex nonword 
condition, where the presence of a real affix drives differential 
parsing efficiency across OS levels; LP indexes the depth and precision 
of lexical-semantic knowledge, including knowledge of stems as free 
morphemes with discrete lexical entries, and its effect on early 
parsing is expressed specifically in the simple nonword condition, 
where stem recognition triggers a parse that cannot be resolved, 
producing a cost that grows with LP. The concurrent modulation of the 
N250 complexity effect by both OS and LP is consistent with 
\citet{beyersmannEmbeddedStemPriming2016}'s suggestion that 
affix-chunking and embedded stem activation mechanisms may operate 
concurrently during early morphological segmentation, with the present 
data extending this proposal by showing that these two mechanisms are 
dissociable across individual difference dimensions: OS indexes the 
efficiency of affix-based parsing, while LP indexes the propensity for 
stem-based parsing independent of affix status.

This mechanistic distinction between OS and LP at the N250 for nonwords 
is further grounded by the word data reported above. For familiar 
words, OS modulated the base frequency effect at the N250 --- a finding 
that reflects OS operating at the level of stem-form entrenchment in 
stored lexical representations. The specificity of the affix-led 
account for nonwords therefore does not imply that OS is exclusively 
sensitive to affixes in all contexts; rather, it suggests that when 
readers encounter a novel form for which no stored whole-word 
representation is available, OS-driven parsing activity is channelled 
specifically through affix recognition, because it is the affix that 
provides the primary morphological signal in that context. LP, by 
contrast, operates through stem recognition in both familiar and novel 
form processing, consistent with its profile as an index of the depth 
and precision of lexical-semantic representations in which stems are 
the primary meaning-bearing units.